\begin{claim}{D.3.5}{Proved}{Cardinality bounds on finite-dimensional targets give the sphere test}
Let $\mathcal F_\kappa$ be the family from D.3.4, where
$\kappa$ is any infinite cardinal. Write
\[
 \mathcal P_2=\{X:\widetilde H_*(X;\mathbb Z[1/2])=0\}.
\]
Then
\[
 \begin{aligned}
 \mathcal N(\mathcal F_\kappa)\cap\mathcal P_2
   &=\Nul\cap\mathcal P_2\\
   &=\{X\in\mathcal P_2:
       \Map_*(X,Z)\simeq *\text{ for every simply connected}\\[-2pt]
   &\hspace{42mm}\text{finite-dimensional CW complex }Z\}.
 \end{aligned}
 \tag{D.3.5.1}
\]
In particular every localizer constructed from the families
$\mathcal F_\kappa$ in D.3.4 is precisely
$W_S\cap R(\loc(\mathbb S/2))$.
Increasing the cardinal bound produces no new element of that
interval. This equality does not decide whether this common
element differs from $W_{\mathrm{rev}}$.
\end{claim}
\begin{proof}
We first prove a countable-source extension of the sphere
criterion. Suppose $X$ has countably many cells and
\[
             \Map_*(X,S^m)\simeq *\quad(m\geq2).
 \tag{D.3.5.2}
\]
Choose an increasing exhaustion $X=\bigcup_{i\geq0}X_i$
by finite pointed subcomplexes. These are finite cell
subcomplexes, not necessarily skeleta of $X$.
For each $a\geq0$ and $m\geq2$ put
\[
       G_i(a,m)=[\Sigma^aX_i,\Omega S^m]_*.
\]
These are countable groups: the source is finite and the
homotopy groups of a sphere are countable, so the finite
obstruction calculation of the mapping set is countable.
The exact sequence for maps out of a sequential homotopy
colimit, applied to $\Sigma^aX$ and $\Sigma^{a+1}X$, gives
\[
       \lim_iG_i(a,m)=0,\qquad
       \lim_i^1G_i(a,m)=*.
\]
By \cite[Proposition~4(b)]{I_Strom2012}, this tower of
countable groups is Mittag--Leffler because its $\lim^1$ vanishes.
Its stable images are the images of its inverse limit:
after passing to stable images the bonding maps are
surjective, so every element of a stable image extends to a
compatible sequence. Since the limit is zero, every stable
image is zero. Thus each tower is \emph{pro-zero}:
for every $i,a,m$ there is $j\geq i$ for which
$G_j(a,m)\to G_i(a,m)$ is zero.

Let $W=\bigvee_{\alpha\in I}S^{n_\alpha}$ be any simply
connected wedge of spheres, of any cardinality. The
Hilton--Milnor decomposition identifies $\Omega W$ with
the weak product of spaces $\Omega S^{d(w)}$, indexed by
basic words $w$ in the sphere labels. This is a weak
product, not the categorical product. For a finite source
$\Sigma^aX_i$, maps and homotopies factor through finite
subproducts, giving a natural bijection of pointed sets
\[
 [\Sigma^aX_i,\Omega W]_*
   \cong\prod_w' [\Sigma^aX_i,\Omega S^{d(w)}]_*.
 \tag{D.3.5.3}
\]
One may verify the assertion first on finite subwedges by
Hilton--Milnor and then pass to the filtered union. Compactness
of the finite source, and of the source of a homotopy, gives
the displayed assertion for the entire wedge. The same
weak-product decomposition is the one used in
\cite[proof of Proposition~6]{I_Strom2012}.

For fixed $i,a$, only finitely many \emph{sphere dimensions}
$d(w)$ can contribute to the right-hand side of
(D.3.5.3): if $d(w)>\dim(\Sigma^aX_i)+1$, the factor is
trivial by connectivity. Choose $j$ large enough that all
the finitely many relevant sphere towers have zero
restriction to stage $i$. This one $j$ works for every
labelled copy of each dimension, regardless of the size of
$I$. Naturality of (D.3.5.3) therefore gives
\[
 [\Sigma^aX_j,\Omega W]_*
       \longrightarrow[\Sigma^aX_i,\Omega W]_*
       \quad\text{equal to zero}.
\]
The tower is pro-zero for every $a$. The sequential
homotopy-colimit exact sequence now gives
$\Map_*(X,\Omega W)\simeq *$ and, in particular,
\[
                \lim_i^1[\Sigma X_i,W]_* = *.
 \tag{D.3.5.4}
\]

It remains to check the unsuspended components of
$\Map_*(X,W)$; loop-space nullity alone would not do so.
By \cite[Corollary~12(a)]{I_Strom2012}, condition
(D.3.5.2) makes every finite subbouquet $V\subset W$
$X$-null. Given $f:X\to W$ and a finite $X_i$, its image
lies in some such $V$. The composite of $f$ with the
retraction $W\to V$ is nullhomotopic. Its restriction to
$X_i$, followed by $V\to W$, is $f|_{X_i}$, so the latter
is nullhomotopic. Thus $f$ represents the zero element of
$\lim_i[X_i,W]_*$, and (D.3.5.4) makes $f$ nullhomotopic.
Apply the same argument to every suspension of $X$.
It follows that $\Map_*(X,W)\simeq *$ for every simply
connected wedge of spheres.

The class $\{Z:\Map_*(X,Z)\simeq *\}$ is a resolving
kernel. \cite[Corollary~12(b)]{I_Strom2012} now puts in
it every simply connected finite-dimensional space.
This proves the countable-source assertion without a
finite-type hypothesis on $X$.

\medskip\noindent
\textbf{Comparison of integral and $2$-local sphere tests.}
For every $X\in\mathcal P_2$ and every simply connected
nilpotent space $Y$, the natural map induces
\[
        \Map_*(X,Y)\simeq\Map_*(X,Y_{(2)}).
 \tag{D.3.5.5}
\]
Here is an arithmetic-square verification, including the
possible infinite size of $X$. Such an $X$ is connected,
rationally acyclic, and mod-$q$ acyclic for every odd
prime $q$. Apply the general arithmetic fracture theorem
\cite[Theorem~13.1.4]{FGHMayPonto2012} first with the set
of all primes and then with the singleton set $\{2\}$.
The completion vertex in the first square is
$\prod_qY_q^{\wedge}$. Each factor for $q\ne2$ has
contractible pointed mapping space from $X$, since it is
mod-$q$-homology-local by
\cite[Theorems~11.1.1--11.1.2]{FGHMayPonto2012}.
Both rational vertices have contractible mapping space
from $X$, since rational nilpotent spaces are rational
homology-local. Mapping into the first square consequently
gives $\Map_*(X,Y)\simeq\Map_*(X,Y_2^{\wedge})$.
The second square gives the same equivalence with
$Y_{(2)}$ in place of $Y$. Naturality proves (D.3.5.5).
These fracture theorems require no finite-generation
hypothesis. Only their special case $Y=S^m$ is needed here.

\medskip\noindent
\textbf{Reduction of arbitrary sources to countable ones.}
Suppose now that $X\in\Nul\cap\mathcal P_2$, with no
cardinality assumption. G.7 supplies the countable
enlargement property for $\Nul$. The class $\mathcal P_2$
has the same property inside $X$: for a countable
simplicial subcomplex, adjoin the finite supports of
bounding chains in $X$ for all its reduced
$\mathbb Z[1/2]$-cycles, and iterate. There are only
countably many such finite chains, since
$\mathbb Z[1/2]$ is countable. Every cycle in the resulting
union bounds at the next stage. This also applies in
reduced degree zero and ensures connectedness.

Alternate the two countable enlargement constructions.
Each of the two null classes is closed under pointed
homotopy colimits, so the sequential union lies in both
classes. Hence every countable collection of simplices
of $X$ lies in a countable subcomplex $X'$ with
$X'\in\Nul\cap\mathcal P_2$. These subcomplexes form
a filtered system with homotopy colimit $X$. The use of
simplicial models and transfer to CW homotopy types is
exactly as in G.7; it makes no assumption that the
singular complex of a countable CW complex is countable.

By (D.3.5.5), each $X'$ satisfies the integral sphere
condition (D.3.5.2). The countable-source assertion
therefore makes $X'$ null into every simply connected
finite-dimensional target. Mapping the filtered homotopy
colimit into a fixed target gives a homotopy limit of
contractible spaces. Thus the same nullity holds for $X$.

Conversely, $\mathcal F_\kappa$ contains every simply
connected sphere. If $X\in\mathcal P_2$ is
$\mathcal F_\kappa$-null, (D.3.5.5) makes it null into
every $2$-local sphere, hence $X\in\Nul$. This proves
all equalities in (D.3.5.1). Finally the aspherical class
of $R(\loc(\mathbb S/2))$ is $\mathcal P_2$, and the
reversible cofiber criterion turns equality of the
displayed null classes into the stated equality of
fundamental localizers.
\end{proof}
